% ===== PRÉAMBULE CANONIQUE — à coller VERBATIM en tête de CHAQUE .tex =====
\documentclass[11pt,a4paper]{article}
\usepackage[utf8]{inputenc}\usepackage[T1]{fontenc}\usepackage{lmodern}
\usepackage[french]{babel}
\usepackage{amsmath,amssymb,mathtools}\usepackage{icomma}
\usepackage[version=4]{mhchem}          % \ce{...} : équations & formules
\usepackage{siunitx}\sisetup{locale=FR} % \qty{}{}, \si{}, \ang{}
\usepackage{chemfig}                    % \chemfig{...} : structures développées
\usepackage{xcolor}\usepackage{colortbl}
\usepackage{tikz}\usepackage{pgfplots}\pgfplotsset{compat=1.18}
\usepackage[most]{tcolorbox}\usepackage{enumitem}\usepackage{array,booktabs}
\usepackage[a4paper,margin=2.2cm]{geometry}
\usetikzlibrary{arrows.meta,calc,angles,quotes,positioning,patterns,backgrounds,shapes.geometric}
\definecolor{primary}{RGB}{222,49,99}\definecolor{primarydark}{RGB}{150,22,55}
\definecolor{defcol}{RGB}{0,114,178}\definecolor{methcol}{RGB}{52,92,148}
\definecolor{excol}{RGB}{0,135,90}\definecolor{attncol}{RGB}{200,40,55}
\tcbset{boxbase/.style={enhanced,breakable,boxrule=0.9pt,arc=2.5pt,
  left=3mm,right=3mm,top=2.3mm,bottom=2.3mm,fonttitle=\bfseries,coltitle=white}}
% --- boîtes TUTORIEL (noms EXACTS, ne pas renommer) ---
\newtcolorbox{cle}[1][]{boxbase,colback=defcol!6,colframe=defcol,
  title={\textsc{À retenir}\ifx\relax#1\relax\relax\else\textmd{ : \itshape #1}\fi}}
\newtcolorbox{methode}[1][]{boxbase,colback=methcol!7,colframe=methcol,sharp corners,
  title={\textsc{Méthode}\ifx\relax#1\relax\relax\else\textmd{ --- \itshape #1}\fi}}
\newtcolorbox{exemplebox}[1][]{boxbase,colback=excol!5,colframe=excol,
  title={\textsc{Exemple}\ifx\relax#1\relax\relax\else\textmd{ : \itshape #1}\fi}}
\newtcolorbox{piege}[1][]{boxbase,colback=attncol!6,colframe=attncol,
  title={\textsc{Piège}\ifx\relax#1\relax\relax\else\textmd{ : \itshape #1}\fi}}
\newtcolorbox{remarque}[1][]{boxbase,colback=black!4,colframe=black!45,
  title={\textsc{Remarque}\ifx\relax#1\relax\relax\else\textmd{ : \itshape #1}\fi}}
% --- boîtes EXERCICE / CORRIGÉ (noms EXACTS, auto-comptées) ---
\newtcolorbox[auto counter]{exo}[1][]{boxbase,colback=primary!4,colframe=primary,
  title={\textsc{Exercice}~\thetcbcounter\ifx\relax#1\relax\relax\else\textmd{ --- \itshape #1}\fi}}
\newtcolorbox[auto counter]{corrige}[1][]{boxbase,colback=excol!4,colframe=excol,
  title={\textsc{Corrigé}~\thetcbcounter\ifx\relax#1\relax\relax\else\textmd{ --- \itshape #1}\fi}}
\newcommand{\R}{\mathbb{R}}\newcommand{\N}{\mathbb{N}}\newcommand{\Z}{\mathbb{Z}}\newcommand{\ds}{\displaystyle}
\begin{document}
\begin{exo}[Les fonctions de l'aspirine]
La structure de l'aspirine (acide acétylsalicylique) est :
\begin{center}
\chemfig{*6(-=-(-(-[:30]OH)=[:90]O)=-(-O-[:150](-[:-150]CH_3)=[:90]O)=-=)}
\end{center}
\begin{enumerate}[label=\alph*)]
  \item Liste \emph{toutes} les fonctions oxygénées présentes.
  \item Y a-t-il une fonction \textbf{éther} dans l'aspirine ? Justifie soigneusement.
  \item Combien de doubles liaisons $\ce{C=C}$ (hors carbonyles) la molécule comporte-t-elle ?
\end{enumerate}
\end{exo}

\begin{exo}[Compter et classer les alcools ; effet d'une réduction]
On considère la molécule \ce{(CH3)2C(OH)-CH2-CH(OH)-CO-CH3}.
\begin{enumerate}[label=\alph*)]
  \item Combien de fonctions alcool porte-t-elle ? Donne la classe de chacune.
  \item Quelle autre fonction contient-elle ?
  \item On réduit le carbonyle ($\ce{C=O}\rightarrow\ce{CH-OH}$). Combien de fonctions alcool la molécule compte-t-elle alors au total, et de quelles classes ?
\end{enumerate}
\end{exo}

\begin{exo}[Azote : amine ou amide ?]
On donne la molécule \ce{H2N-CH2-CO-NH-CH2-COOH} (un dipeptide, la glycylglycine).
\begin{enumerate}[label=\alph*)]
  \item Identifie les \textbf{trois} fonctions présentes.
  \item La molécule contient deux atomes d'azote : portent-ils la même fonction ? Justifie.
  \item L'enchaînement central $-\ce{CO-NH}-$ est-il une amine secondaire ? Explique l'erreur classique.
\end{enumerate}
\end{exo}

\begin{exo}[Une molécule aromatique polyfonctionnelle]
Un composé aromatique a la structure suivante :
\begin{center}
\chemfig{*6(-(-CHO)=-(-OCH_3)=(-CH_2OH)-=)}
\end{center}
\begin{enumerate}[label=\alph*)]
  \item Liste les fonctions présentes (celles du cycle et de ses substituants).
  \item Le substituant $-\ce{OCH3}$ porté par le cycle est-il une fonction \textbf{éther} ou \textbf{ester} ? Justifie.
  \item Combien de doubles liaisons $\ce{C=C}$ le cycle contient-il ?
\end{enumerate}
\end{exo}

\begin{exo}[Vrai / faux sur une molécule polyfonctionnelle]
Soit la molécule \ce{CH2=CH-CH2-CO-CH2-CHO}. Réponds par \textbf{vrai} ou \textbf{faux} en justifiant.
\begin{enumerate}[label=\alph*)]
  \item Elle contient deux fonctions carbonyle.
  \item Elle contient une fonction acide carboxylique.
  \item Elle contient trois fonctions chimiques.
  \item Elle contient une fonction alcool.
  \item L'un de ses carbonyles est un aldéhyde situé en bout de chaîne.
\end{enumerate}
\end{exo}

\end{document}
